%----------------------------------------------------------------------------- %%
\chapter{\babDua}
\label{bab:2}
%----------------------------------------------------------------------------- %%

This chapter reviews the theoretical concepts and prior work that form the foundation of this study. Section 2.1 outlines the Kubernetes architecture and its resource management principles. Section 2.2 defines the resource fragmentation problem within dynamic clusters. Section 2.3 details resource defragmentation and multi-resource bin packing. Section 2.4 describes the Kubernetes Descheduler and its current limitations. Section 2.5 explains predicted-target eviction-candidate selection. Section 2.6 discusses runtime load-aware eviction guarding. Section 2.7 discusses communication-aware and network-aware scheduling approaches. This chapter concludes with Section 2.8, which identifies the current approaches and the resulting research gap.

\section{Kubernetes Architecture and Resource Management}
\label{sec:kubernetesArchitecture}

Kubernetes distributes Pods across worker Nodes. For CPU and memory, a container may declare a resource request, a resource limit, both, or neither. Requests represent the resource baseline used by the scheduler when determining whether a Pod fits on a Node, whereas limits are runtime enforcement ceilings applied by the kubelet and container runtime~\citep{k8s_resource_management}. Kubernetes uses effective Pod requests when evaluating placement feasibility and resource allocation on each Node. Resource limits do not participate in this fit calculation; they define runtime enforcement ceilings after the Pod is running. Scheduler scoring strategies such as \textit{MostAllocated} and \textit{RequestedToCapacityRatio} may then use request-based allocation to favor denser packing~\citep{k8s_resource_bin_packing}.

Requests and limits also contribute to Pod Quality of Service (QoS) classification. A Pod is \textit{Guaranteed} when every container declares CPU and memory requests equal to its corresponding limits, while Pods with partial or unequal declarations are generally \textit{Burstable}. A Pod whose containers declare neither CPU nor memory requests or limits is \textit{BestEffort}~\citep{k8s_pod_qos}. If a limit is declared without a request, Kubernetes can use that limit as the effective request for the resource. A namespace \textit{LimitRange} may also assign default requests or limits during Pod admission. These defaults become part of the Pod's effective resource configuration~\citep{k8s_limit_range}.

These decisions are made when a Pod is scheduled. A running Pod is not continuously relocated when the cluster state changes, so an initially acceptable layout can become inefficient after scale-outs, terminations, and replacement events. This separation between initial placement and later cluster state motivates post-placement correction.

\section{Resource Fragmentation and Stranded Capacity}
\label{sec:resourceFragmentation}

Resource fragmentation occurs when free capacity is divided among Nodes in shapes that do not match incoming workloads. A cluster can have enough aggregate CPU and memory while no individual Node satisfies both dimensions of a Pod request. This makes CPU and memory placement a multi-dimensional problem rather than a scalar load-balancing problem.

CPU-memory skew is one source of stranded capacity. If one Node is CPU-heavy and another is memory-heavy, each leaves a different residual resource that cannot be consumed alone. The absolute difference between normalized CPU and memory fractions provides a direct shape indicator, while the product of free CPU and memory fractions represents the jointly usable free-space region. Prior work shows that directed eviction can repair fragmented Kubernetes layouts and support Node reclamation~\citep{larsson_klein_elmroth_2023}.

\section{Resource Defragmentation and Multi-Resource Bin Packing}
\label{sec:binPacking}

Resource packing seeks to place workloads on fewer active Nodes so that empty Nodes can be removed or powered down. In a multi-resource setting, dense packing alone is insufficient: a target can be highly utilized but badly skewed. A useful packing score should therefore represent both density and CPU-memory balance.

The policy in this research uses normalized CPU and memory fractions. Average utilization represents density, while a bin score combines density with the absolute distance between the two dimensions. The Resource Imbalance Index (RII) preserves the direction of skew, and the Free-Space Index (FSI) measures the rectangular CPU-memory headroom. Together these quantities identify Nodes that are lightly loaded, poorly shaped, or under pressure from a narrow free-space region.

\section{Kubernetes Descheduler}
\label{sec:deschedulerLimitations}

The Kubernetes Descheduler evicts Pods that violate a policy or contribute to an undesirable placement state, after which normal controllers recreate the Pods and kube-scheduler places them again~\citep{k8s_descheduler_docs}. It provides an appropriate extension point for node reclamation because it can repair placement without modifying the scheduler.

An eviction policy must nevertheless account for the scheduler boundary. Deleting a Pod does not guarantee that it moves to a useful Node; it may become Pending, return to its origin, or land on a target that worsens fragmentation. A node reclamation-oriented policy therefore benefits from predicting feasible non-origin targets and rejecting movements that exceed a configured ceiling or pack downward.

Upstream utilization policies provide important baselines. \textit{HighNodeUtilization} packs by draining low-utilization Nodes toward more utilized Nodes, but it does not explicitly optimize CPU-memory stranding. This distinction is central to the evaluation in Chapter~\ref{bab:4}.

\section{Predicted-Target Eviction-Candidate Selection}
\label{sec:predictedTargetSelection}

Eviction-candidate selection can be formulated around the expected landing state. For every evictable Pod on a drain candidate, \textit{predictSchedulerTarget} simulates request-based placement on non-origin Nodes that satisfy scheduler constraints. Feasible targets must have sufficient CPU and memory, remain below a node reclamation ceiling after placement, and be at least as packed as the source.

The target with the best projected CPU-memory bin score is retained for each Pod. The Pod whose predicted target has the highest score is selected for eviction. This single criterion aligns the action directly with the desired destination state: the selected movement should place complementary resource shapes together while preserving node reclamation pressure.

\section{Runtime Load-Aware Eviction Guard}
\label{sec:runtimeLoadAwareEviction}

Request-based resource information is sufficient for scheduler feasibility and bin-packing analysis because Kubernetes places Pods according to their effective requests rather than their instantaneous consumption. However, the operational cost of eviction depends on the Pod's runtime state. A Pod that is attractive from a defragmentation perspective may still be handling high live demand when the Descheduler attempts to evict it. Load-aware orchestration addresses this limitation by incorporating runtime resource usage into placement decisions. For example, \citet{marchese2024loadaware} propose a Kubernetes orchestration strategy that continuously considers dynamic resource usage patterns when placing microservice applications in shared clusters.

This research applies the same principle at the eviction boundary rather than replacing the scheduler. The defragmentation policy evaluates source Nodes, target feasibility, and projected CPU-memory shape in request space because those quantities determine whether the replacement Pod can be scheduled. The actual-usage guard then compares live consumption with the same request baseline to determine whether the selected Pod is currently active relative to the capacity used for placement. Resource limits are not used as denominators because usage relative to a limit indicates proximity to CPU throttling or a memory out-of-memory boundary. This represents resource saturation rather than the busy-Pod protection objective of this plugin. Limits can nevertheless affect observed usage indirectly by throttling CPU or terminating a container that exceeds its memory limit~\citep{k8s_resource_management}.

This Descheduler guard also has a different objective from kubelet node-pressure eviction. Under node pressure, the kubelet ranks Pods to free a pressured resource based on whether usage exceeds requests, Pod Priority, and usage relative to requests. QoS influences this ranking through the requests and priority characteristics of each Pod~\citep{k8s_node_pressure_eviction}. The Actual Usage Evictor instead decides whether a strategy-selected eviction should proceed without disrupting a currently busy Pod. This separation preserves request-space defragmentation while adding a runtime safety decision immediately before eviction.

\section{Communication-Aware and Network-Aware Placement}
\label{sec:networkAwareScheduling}

Microservices exchange traffic across service dependencies. Placement across distant zones or degraded links can increase latency even when compute resources fit. Communication-aware Kubernetes studies therefore incorporate topology, traffic, or measured network conditions into placement decisions~\citep{marchese2022network,marchese_tomarchio_2022,chen2024trade,pusztai_pogonip_2021}.

This research applies network awareness as an eviction guardrail. A global pre-eviction filter compares the current communication cost with feasible alternatives, while strategy-level filtering can evaluate a narrower target set. The design keeps resource improvement and network protection separable but composable.

\section{Current Approaches and Research Gap}
\label{sec:researchGap}

Existing mechanisms address individual parts of the problem. Scheduler scoring improves new placement, utilization-based descheduling drains Nodes according to thresholds, fragmentation research motivates directed eviction, load-aware orchestration highlights the importance of runtime demand, and network-aware schedulers protect communication locality. The remaining gap is a practical Descheduler pipeline that combines:

\begin{itemize}
    \item source assessment using utilization and CPU-memory shape;
    \item RII/FSI-based drain ordering;
    \item request-feasible scheduler-target prediction;
    \item a direct projected-target balance criterion;
    \item bounded, convergent eviction;
    \item actual-usage pre-eviction filtering for busy Pods; and
    \item network-aware filtering for latency-sensitive workloads.
\end{itemize}

\textit{ResourceDefragmentation} and the Network-Aware Evictor are designed to fill this gap while leaving final binding to the standard kube-scheduler.
